The route from the transition matrix to slow ACF decay runs through the spectrum of $\mathbf T$. Let $\mathbf v_k$ and $\mathbf w_k$ denote the right and left eigenvectors of $\mathbf T$ at non-unit eigenvalue $\lambda_k$, normalised biorthonormally ($\mathbf w_k^\top \mathbf v_l = \delta_{kl}$), and let $\sigma_{|G|}^2$ denote the stationary variance of $|G_t|$ under the marginal mixture~\eqref{eq:mixture}. The closed-form bilinear identity for the absolute growth-rate ACF of a stationary CHMM, $\E[|G_t|\,|G_{t+\tau}|] = \mathbf m^\top \mathrm{diag}(\bar{\boldsymbol\pi})\,\mathbf T^\tau\,\mathbf m$ for lags $\tau \ge 1$, with $\mathbf m_k = \E[|G_t| \mid s_t = k]$, and its eigendecomposition over the non-unit eigenvalues of $\mathbf T$ are standard in the regime-switching literature~\cite{hamilton1994time, timmermann2000moments}; at $\tau = 0$ the identity does not apply, since the conditional second moment replaces $m_k^2$. Under irreducibility and aperiodicity of $\mathbf T$ (with $\bar{\boldsymbol\pi}$ its unique stationary distribution), finite per-state second moments, and diagonalizability of $\mathbf T$, the spectral decomposition $\mathbf T^\tau = \mathbf 1\bar{\boldsymbol\pi}^\top + \sum_{k=2}^K \lambda_k^\tau\, \mathbf v_k \mathbf w_k^\top$ subtracts the marginal product to give, for $\tau \ge 1$,
\begin{equation}
    \rho_{|G|}(\tau) = \sum_{k=2}^{K} a_k\, \lambda_k^{\tau},
    \qquad
    a_k = \frac{\big(\mathbf m^\top \mathrm{diag}(\bar{\boldsymbol\pi})\, \mathbf v_k\big)\big(\mathbf w_k^\top \mathbf m\big)}{\sigma_{|G|}^2}.
    \label{eq:acf_normalised}
\end{equation}
This is a real-valued sum of geometric modes at the non-unit eigenvalues of $\mathbf T$; complex-conjugate pairs combine into damped oscillatory modes, and a non-diagonalisable $\mathbf T$ yields polynomial-times-geometric Jordan terms instead. Wherever a dominant-mode share is reported below, it is defined as $\max_k |a_k \lambda_k| / \sum_{k \ge 2} |a_k \lambda_k|$ at lag $1$: a share of the total \emph{absolute} modal contribution, which coincides with a share of $\rho_{|G|}(1)$ itself only when all contributions carry the same sign. Our contribution is purely empirical. We apply equation~\eqref{eq:acf_normalised} to recast the Ryd\'{e}n et al.~\cite{ryden1998stylized} low-$K$ failure on the SPY 2014--2024 instance as a rank statement on $\mathbf T - \mathbf 1\bar{\boldsymbol\pi}^\top$. At $K = 2$, the ACF reduces to the mono-exponential $\rho_{|G|}(\tau) = a_2\lambda_2^\tau$ with $\lambda_2 = T_{11} + T_{22} - 1$: its persistence can be tuned through $T_{11} + T_{22}$ but its shape is restricted to one geometric mode. The fitted SPY $K = 2$ chain confirms this empirically: the single non-unit eigenvalue trivially carries $100\%$ of the lag-1 modal contribution, and it is strongly persistent ($\lambda_2 = 0.94$ for Gaussian emissions, $0.95$ for the shared-$\nu$ Student-$t$), so the slow decay is already present at the lower bound. At $K \ge 3$ the matrix admits multiple non-unit eigenvalues and the marginal mixture admits more components. The empirical results test whether those finite-mode and finite-mixture restrictions are active on the studied data; they do not imply that a finite-state HMM reproduces general power-law or multi-scale decay.
